\documentclass[11pt]{article}
\usepackage[margin=1in]{geometry}
\usepackage{amsmath,amssymb,amsthm}
\usepackage{hyperref}
\usepackage{graphicx}

\newtheorem{theorem}{Theorem}
\newtheorem{corollary}{Corollary}
\newtheorem{remark}{Remark}

\newcommand{\B}{B}
\newcommand{\Sph}{S}
\newcommand{\eps}{\varepsilon}
\newcommand{\weakstar}{w^*}

\title{Strong Exposure as a Sufficient Condition for Bidual Continuity}
\author{Run fa\_banach\_001}
\date{2026-07-01}

\begin{document}
\maketitle

\section*{Source and Target}

The source paper is S. Daptari, V. Montesinos, and T. S. S. R. K. Rao,
\emph{A study of weak$^*$-weak points of continuity in the unit ball of dual spaces},
arXiv:2506.08458.

In Section 5.2, after a higher-dual preservation proposition, the authors write:

\begin{quote}
Related to Lemma 5.1, let us mention that we don't know of any geometric condition on
$x\in S(X)$ in terms of $X$ alone for ensuring that $x$ should be a point of continuity of
$id:(B(X^{**}),w^*)\to (B(X^{**}),w)$.
\end{quote}

\begin{center}
\includegraphics[width=0.92\linewidth]{figures/open_problem_crop.png}
\end{center}

\section*{Result}

We give a direct sufficient condition stated entirely in the original space.
The conclusion is stronger than requested: it gives weak-star-to-norm continuity.

\begin{theorem}
Let $X$ be a real or complex Banach space. Suppose that $x\in S_X$ is a strongly exposed
point of $B_X$, strongly exposed by $x^*\in S_{X^*}$, so that
$\operatorname{Re}x^*(x)=1$ and, for every $\eps>0$, there is $\delta>0$ such that
\[
 y\in B_X,\quad \operatorname{Re}x^*(y)>1-\delta
 \quad\Longrightarrow\quad
 \|y-x\|<\eps .
\]
Then the canonical image of $x$ in $X^{**}$ is a point of continuity of
\[
id:(B_{X^{**}},\sigma(X^{**},X^*))\longrightarrow (B_{X^{**}},\|\cdot\|).
\]
In particular, $x$ is a point of continuity of
$id:(B_{X^{**}},w^*)\to (B_{X^{**}},w)$.
\end{theorem}

\section*{Proof Intuition}

Strong exposure says that one weak slice cut out by $x^*$ gives arbitrarily small
norm neighborhoods of $x$ inside the original unit ball. Goldstine's theorem says that
$B_X$ is weak-star dense in $B_{X^{**}}$. Therefore a bidual point lying in the same
weak-star slice is a weak-star limit of original-space points eventually lying in the
small norm slice. The norm ball around $x$ is weak-star closed in the bidual, so the
bidual point lies in the same small norm ball.

\section*{Proof}

Fix $\eps>0$. By strong exposure, choose $\delta>0$ such that
\[
 y\in B_X,\quad \operatorname{Re}x^*(y)>1-\delta
 \quad\Longrightarrow\quad
 \|y-x\|<\eps .
\]
We claim that the weak-star slice
\[
 U_\delta=\{u\in B_{X^{**}}:\operatorname{Re}u(x^*)>1-\delta\}
\]
is contained in $x+\eps B_{X^{**}}$.

Let $u\in U_\delta$. By Goldstine's theorem, there is a net $(y_\alpha)$ in $B_X$
which converges to $u$ in $\sigma(X^{**},X^*)$. Since
$\operatorname{Re}u(x^*)>1-\delta$, eventually
$\operatorname{Re}x^*(y_\alpha)>1-\delta$. Hence eventually
$\|y_\alpha-x\|<\eps$, i.e. $y_\alpha\in x+\eps B_X\subset x+\eps B_{X^{**}}$.
The set $x+\eps B_{X^{**}}$ is weak-star closed, because $B_{X^{**}}$ is weak-star
compact and translations are weak-star homeomorphisms. Therefore
$u\in x+\eps B_{X^{**}}$.

Thus every norm neighborhood of $x$ in $B_{X^{**}}$ contains a weak-star neighborhood
of $x$ relative to $B_{X^{**}}$. This is exactly weak-star-to-norm continuity at $x$.
Since the weak topology on $X^{**}$ is weaker than the norm topology, the asserted
weak-star-to-weak continuity follows.

\section*{Verification and Limitations}

The proof is purely functional analytic and uses only the definition of strong exposure
and Goldstine's theorem. It gives a sufficient condition, not a characterization. It
also says nothing about the later question in the source paper concerning stability of
the dense $G_\delta$ property under subspaces, quotients, or $c_0$-sums.

Run-level novelty checks were bounded: the lightweight indexes were searched for
arXiv:2506.08458 and the main phrases ``strongly exposed'', ``bidual'', ``points of
continuity'', and ``weak-star weak''. No prior run packet for this result was found.
A local source search suggests the underlying lift of strong exposure to the bidual is
standard, so the recommendation is to treat this as a useful identified sufficient
condition rather than as a high-novelty theorem.

\section*{References}

\begin{itemize}
\item S. Daptari, V. Montesinos, and T. S. S. R. K. Rao,
\emph{A study of weak$^*$-weak points of continuity in the unit ball of dual spaces},
arXiv:2506.08458.
\item J. Diestel, \emph{Geometry of Banach Spaces: Selected Topics}, Lecture Notes in
Mathematics 485, Springer, 1975. Background on denting and exposed points.
\end{itemize}

\end{document}

